\documentclass[journal]{IEEEtran}

% --- PAQUETES BÁSICOS ---
\usepackage[utf8]{inputenc}
\usepackage[T1]{fontenc}
\usepackage{amsmath, amsfonts, amssymb}
\usepackage{graphicx}
\usepackage[siunitx]{circuitikz} 

%Traducciónes específicas
\renewcommand{\abstractname}{Resumen}
\renewcommand{\tablename}{Tabla}
\renewcommand{\figurename}{Figura}
\renewcommand{\IEEEkeywordsname}{Palabras Clave}

\begin{document}

%Título y autores
\title{Informe 1: Circuitos Resistivos}
\author{Maria José Álvarez Hernández, \textit{123456789}, \textit{Universidad de Antioquia}, 
        y Cristian Alejandro Guarin Marin, \textit{1041228448}, \textit{Universidad de Antioquia}}

\maketitle

%Resumen
\begin{abstract}
En este documento se detalla la comprobación práctica de la ley de Ohm, las leyes de Kirchhoff y las técnicas de simplificación de circuitos resistivos. 
Se analizaron tres configuraciones principales: paralela, serie y mixta (requiriendo transformación delta-estrella). La metodología consistió en contrastar 
los valores teóricos calculados en el preinforme, los resultados obtenidos mediante simulación en LtSpice y las mediciones físicas tomadas en el laboratorio. 
A través de esta comparación, se evaluó la precisión de los métodos teóricos y se analizó el impacto de las tolerancias de los componentes reales.
\end{abstract}

%Palabras clave
\begin{IEEEkeywords}
Ley de Kirchhoff, Ley de Ohm, Multímetro, Protoboard, Circuitos Resistivos.
\end{IEEEkeywords}

%Introducción
\section{Introducción}


%Metodología ------------------------------------------------------------------------------------------------------------------------------------------------------------------
\section{Metodología}
Para el desarrollo experimental se empleó una placa de pruebas (protoboard), resistores de diferentes valores comerciales, un multímetro digital y una 
fuente de alimentación de corriente continua configurada a $8V$. 

El procedimiento se dividió en las siguientes fases:
\begin{itemize}
    
    \item \textbf{Selección y medición de componentes:} Inicialmente, se midió el valor real de resistencia de cada componente con el multímetro y se comparó 
    con el valor teórico dictado por su código de colores. \textit{Cabe aclarar que, para el desarrollo de este montaje físico, se decidió utilizar valores de 
    resistencia distintos a los propuestos inicialmente en el preinforme. Esto se hizo con el propósito de evitar que el circuito resultara demasiado simple o 
    redundante, garantizando así un escenario de medición que permitiera observar y analizar los fenómenos eléctricos de manera más rigurosa.}
    
    \item \textbf{Circuito en paralelo (Figura 1):} Se realizó el montaje de resistores en paralelo. Con la fuente desconectada, se midió la resistencia 
    equivalente del sistema. Posteriormente, se activó la fuente y se midieron tanto la corriente total suministrada como las corrientes de rama individuales.
    
    \item \textbf{Circuito en serie (Figura 2):} Se implementó una configuración en serie, midiendo la resistencia equivalente total y las caídas de tensión (voltaje) en 
    cada uno de los resistores que componen el circuito.
    
    \item \textbf{Circuito mixto (Figura 3):} Se armó una topología que requiere una transformación delta-estrella para su simplificación teórica. 
    En este montaje se midió experimentalmente la resistencia equivalente vista desde la fuente y la corriente total suministrada por la misma.

\end{itemize}
En todos los casos, los circuitos fueron previamente simulados y resueltos de forma manual para obtener los puntos de operación teóricos.

%Resultados ------------------------------------------------------------------------------------------------------------------------------------------------------------------
\section{Resultados}
A continuación se presentan los datos recolectados durante la práctica. En la Tabla \ref{tab:resistencias} se detallan los valores nominales y medidos de los 
resistores utilizados. Posteriormente, se exponen las mediciones de corriente, voltaje y resistencia equivalente para cada una de las configuraciones, contrastando los 
datos del preinforme, la simulación y el montaje físico.

%Tabla de resistencias ------------------------------------------------------------------------------------------------------------------------------------------------------------------
\begin{table}[htbp]
\centering
\caption{Valores de las Resistencias}
\label{tab:resistencias}
\begin{tabular}{|c|c|c|}
\hline
\textbf{Resistencia} & \textbf{Valor estimado} & \textbf{Valor medido}  \\
\hline
$R_a$ & $1500\varOmega$ & $1482\varOmega$ \\
\hline 
$R_b$ & $2200\varOmega$ & $2212\varOmega$ \\
\hline 
$R_c$ & $3300\varOmega$ & $3194\varOmega$ \\
\hline
$R_d$ & $1000\varOmega$ & $989\varOmega$ \\
\hline
$R_e$ & $2000\varOmega$ & $1948\varOmega$ \\
\hline
\end{tabular}
\end{table}

\subsection{Circuito 3}

%Circuito 3 -----------------------------------------------------------------------------------------------------------------------------------------------------------------
\begin{figure}[htbp]
    \centering
    \begin{circuitikz}
        
        %Fuente alineada a la izquierda
        \draw (0,0) to[battery1, invert, l=$V_s$] (0,7);
        
        %Riel superior
        \draw (0,7) to[short, i=$i$] (2,7);
        \draw (2,7) to[short] (7,7);
        
        %Riel inferior
        \draw (0,0) to[short] (7,0);

        %Primer paralelo
        \draw (2,7) to[R, l=$R_a$] (2,3.5);
        \draw (2,3.5) to[R, l=$R_d$] (2,0);

        %Horizontal
        \draw (2,3.5) to[R, l=$R_c$] (7,3.5);
        %Segundo paralelo
        \draw (7,7) to[R, l=$R_b$] (7,3.5);
        \draw (7,3.5) to[R, l=$R_e$] (7,0);
    
    \end{circuitikz}
    
    \caption{Esquemático del Circuito 3.}
    \label{fig:circuito3}
\end{figure}

%Tabla
\begin{table}[htbp]

\centering
\caption{Resultados del Circuito 3}
    
        \begin{tabular}{|c|c|c|c|}
            
            \hline
            & \textbf{Cálculos} & \textbf{Simulación} & \textbf{Montaje} \\
            \hline
            $\mathbf{R_{eq}}$ &  &  & $1549\varOmega$\\
            \hline
            $\mathbf{I}$ &  & $5.11mA$ & $5.15mA$\\
            \hline

    \end{tabular}

\end{table}

%Análisis de resultados------------------------------------------------------------------------------------------------------------------------------------------------------------------
\section{Análisis de Resultados}

%Conclusiones-------------------------------------------------------------------------------------------------------------------------------------------------------------------
\section{Conclusiones}
A partir del desarrollo de la práctica, se logró comprobar de manera experimental la validez de la ley de Ohm y las leyes de Kirchhoff. 
Se verificó que, en los circuitos en paralelo, la corriente total se divide entre las distintas ramas, mientras que en los circuitos en serie, el 
voltaje total se reparte entre los resistores. 

Asimismo, se demostró que las técnicas de simplificación matemática, como la transformación delta-estrella, son herramientas precisas para predecir el comportamiento físico de un circuito. A pesar de los pequeños porcentajes de error inherentes a los componentes físicos, los modelos teóricos y simulados representan fielmente la dinámica del montaje real.

\end{document}